\section{Paper Generation Pipeline}\label{sec:paper}

\subsection{Pipeline overview}\label{sec:paper:overview}

\begin{figure}[t]
  \centering
  \resizebox{\columnwidth}{!}{% F06: paper-generation swim lane — paper-style three-lane workflow.
% Lanes are coloured stage containers; numbered step markers sit to the
% left of each station; the orange refine back-edge is on top.
% style.tikzstyles is loaded by shared/preamble.tex / standalone wrapper.
\begin{tikzpicture}[node distance=1.0cm and 1.6cm]

  % --- six stations laid out in a 3-row grid -------------------------
  \node[station]                                       (i1)  {\figlabel{F6.brainstorm}};
  \node[station, right=2.6cm of i1]                    (i2)  {\figlabel{F6.refine_idea}};

  \node[station emph, below=1.4cm of i1]               (e1)  {\figlabel{F6.bfts}};
  \node[station, right=2.6cm of e1]                    (e2)  {\figlabel{F6.replicate}};

  \node[station, below=1.4cm of e1]                    (w1)  {\figlabel{F6.draft}};
  \node[station emph, right=2.6cm of w1]               (w2)  {\figlabel{F6.review_refine}};

  % --- numbered step markers, perched on the upper-left corner of each
  % station (PaperBench-style badge), so they never overlap any arrow.
  \node[numbered step, above left=0.05cm and -0.2cm of i1.north west] {1};
  \node[numbered step, above left=0.05cm and -0.2cm of i2.north west] {2};
  \node[numbered step, above left=0.05cm and -0.2cm of e1.north west] {3};
  \node[numbered step, above left=0.05cm and -0.2cm of e2.north west] {4};
  \node[numbered step, above left=0.05cm and -0.2cm of w1.north west] {5};
  \node[numbered step, above left=0.05cm and -0.2cm of w2.north west] {6};

  % --- intra-lane forward arrows (thick, paper-style) ----------------
  \draw[big arrow] (i1) -- (i2);
  \draw[big arrow] (e1) -- (e2);
  \draw[big arrow] (w1) -- (w2);

  % --- inter-lane forward arrows: short straight verticals ----------
  \draw[big arrow] (i2.south) -- (e2.north);
  \draw[big arrow] (e2.south) -- (w2.north);

  % --- refine back-edge (orange, on top, label in middle) ----------
  \draw[big arrow, draw=ari-vermilion] (w2.south) to[bend left=35]
       node[edge label, fill=ari-orange!15, font=\sffamily\bfseries\footnotesize, midway]
       {\figlabel{F6.refine}}
       (w1.south);

  % --- lane backgrounds (in background layer) ----------------------
  \begin{scope}[on background layer]
    \node[stage container, fit=(i1)(i2), inner sep=10pt]   (laneI) {};
    \node[stage container, fit=(e1)(e2), inner sep=10pt]   (laneE) {};
    \node[stage container, fit=(w1)(w2), inner sep=10pt]   (laneW) {};
  \end{scope}
  % lane labels OUTSIDE on the right, after the lanes are drawn (top layer)
  \node[stage title, right=4pt of laneI.east, anchor=west] {\figlabel{F6.lane.idea}};
  \node[stage title, right=4pt of laneE.east, anchor=west] {\figlabel{F6.lane.experiment}};
  \node[stage title, right=4pt of laneW.east, anchor=west] {\figlabel{F6.lane.write}};
\end{tikzpicture}
}
  \caption{Paper-generation swim lane: idea, experiment, write. Each
           lane corresponds to a phase; the back edge from
           review to draft is the only non-monotonic transition
           (\cref{sec:paper:loop}).}
  \label{fig:swim}
\end{figure}

Once exploration terminates, the surviving lineage is handed to the
paper-generation pipeline, the second tree of activity in ARI
(\cref{fig:swim}). A \term{write} phase compiles the lineage into a
draft, and a \term{review} phase scores that draft against a rubric;
the two are coupled by a refine loop (\cref{sec:paper:loop}). The
pipeline assembles, from the per-node artefacts left by exploration,
the evidence the writer is allowed to draw on, and persists three
outputs: the \LaTeX{} draft, the per-leaf review with its rationale,
and the rubric instance used to score the run (kept because the
rubric may be generated per paper rather than fixed in advance).

\subsection{Rubric formalism}\label{sec:paper:rubric}

We model the \term{rubric} $\Rubric$ as a finite tree whose leaves
carry tuples
$(\rubricitem{i},\rubricweight{i},\judgefn{i})$. Each leaf $i$ has
a category descriptor $\rubricitem{i}$, a non-negative weight
$\rubricweight{i}$ with $\sum_i \rubricweight{i}=1$, and a
\emph{judge function} $\judgefn{i}:\paper\to[0,1]$ that returns a
graded score. The paper score is the convex combination
\begin{equation}\label{eq:paperscore}
  \paperscore{\paper} \;=\; \sum_i \rubricweight{i}\,\judgefn{i}(\paper).
\end{equation}
This tree-structured rubric is shared with PaperBench, which ARI
reuses through a vendored package (\cref{sec:paper:paperbench}) so
that its grade and aggregation semantics track upstream verbatim.
Two judge families coexist: an LLM-backed judge for qualitative
leaves, and a deterministic judge for leaves with a binary or numeric
verifier (e.g.\ ``is the cost reported in absolute dollars?''). Both
satisfy the same \texttt{Evaluator} interface, so the search engine
(\cref{sec:exploration}) and the paper pipeline reduce axes to a
scalar through one seam.

A rubric is generated in one of two \emph{modes}. In
\emph{agent-benchmark} mode the tree decomposes the paper by
scientific contribution and each leaf asks whether a candidate
submission's executed output reproduces the paper --- the original
PaperBench framing. In \emph{paper-audit} mode the tree's top-level
children are a fixed set of reproducibility axes and each leaf grades
the descriptive completeness of the paper text itself --- a framing
inversion aimed at HPC reproducibility audits, where the question is
not ``can an agent reproduce this paper?'' but ``does this paper
describe enough to be reproducible?''. Venue knowledge (which axes, how
they are phrased) is supplied by a per-venue template rather than
baked into the engine, so ARI core stays domain-agnostic (P4) and a
new venue is a data change, not a code change.

Lifting the paper-audit score out of a degenerate regime --- where an
empty submission drives every submission-oriented leaf to ``no'' ---
required re-phrasing the leaf questions to interrogate the \emph{paper}
rather than a (missing) submission, feeding the judge the paper's
figures alongside its text so vision-capable models can use them, and
admitting optional artifact-description and artifact-evaluation
appendices as extra input. These adjustments live entirely in ARI's
wrapper layer; the vendored judge is untouched, which keeps scores
comparable across upstream versions. We deliberately do \emph{not}
synthesise a reproduction log from the paper's own reported numbers:
doing so would short-circuit the executed-submission stage and inflate
the score in a way not comparable to leaderboard runs.

\subsection{Review/refine loop}\label{sec:paper:loop}

The refine cycle iteratively rewrites $\paper_t$ to $\paper_{t+1}$
under the review feedback $J(\paper_t)$:
\begin{equation}\label{eq:refine}
  \paper_{t+1} \;=\; \refine\!\bigl(\paper_t, J(\paper_t)\bigr).
\end{equation}
The cycle stops when the reviewer accepts the draft --- modelled as the
\term{convergence criterion}
$\paperscore{\paper_{t+1}}-\paperscore{\paper_t}<\epsilon$ --- or when
a small per-section revision cap is reached. Which branch fires more
often is an empirical question \cref{chapter:eval} is meant to answer
once a frozen checkpoint exists; we make no quantitative claim here.

The refine step is intentionally non-monotonic per axis: fixing a
clarity issue may slightly degrade a numeric axis. This is the source
of the back edge in \cref{fig:swim}, and it is why the convergence
criterion is defined on the aggregate score rather than per-axis ---
together with a per-axis floor (\cref{sec:paper:failures}) that
prevents the aggregate from hiding a collapse on any single axis.

\subsection{PaperBench integration}\label{sec:paper:paperbench}

The replication path is \emph{not} a re-implementation of PaperBench:
the upstream package is vendored unchanged so that ARI tracks its task
and rubric semantics verbatim, and ARI contributes only a thin wrapper
around it. The wrapper exposes the protocol-faithful loop as three
composable stages with a shared vocabulary (\cref{fig:pb-stages}):

\begin{enumerate}
  \item \emph{rollout} --- a replication agent works from the paper
        towards a submission (code, and optionally a reproduction
        script);
  \item \emph{reproduce} --- the submission is executed in an isolated
        container, on the local machine or dispatched to a cluster;
  \item \emph{grade} --- the vendored judge scores the submission
        against the rubric tree.
\end{enumerate}

\begin{figure*}[t]
  \centering
  \resizebox{\textwidth}{!}{% F16: the PaperBench replication path as three composable stages, laid out
% as a left-to-right staircase: (1) ROLLOUT turns the paper + rubric into a
% submission (code, reproduce.sh); (2) REPRODUCE executes that submission in
% an isolated sandbox, capturing the reproduction log and outputs; (3) GRADE
% opens up into the weighted rubric tree — the vendored judge assigns each
% leaf a binary grade g (dashed flow), and the paper score is the
% weight-aggregated roll-up s = sum_i w_i g_i (eq:paperscore). Pale-green
% fill marks verdicts (granted leaves, final score); the failed leaf stays
% unfilled; peach marks execution; artifacts and the vendored judge are
% white. House design language comes from _style.tex below; the legacy
% baseline style.tikzstyles is loaded by shared/preamble.tex / the
% standalone wrapper and is not used here.
%
% _style.tex resolves from the build cwd: the en/ja/zh documents \input this
% figure as ../shared/figures/tikz/<fig>.tex (cwd = language dir), the
% render harness compiles from the report root, and the bare name covers
% compiling next to the file itself.
\IfFileExists{../shared/figures/tikz/_style.tex}%
  {% shared/figures/tikz/_style.tex — ARI house design language for system figures.
% ---------------------------------------------------------------------------
% \input this file at the TOP of a figure source, BEFORE \begin{tikzpicture}.
% It only DEFINES colors / \tikzset styles / pics — it typesets nothing — so
% it is safe to \input from several figures in the same document and safe to
% load in any engine (pdflatex / lualatex / xelatex): xcolor + TikZ core only,
% no fontspec, no external icon packages (glyphs below are stroke-drawn pics).
%
% Relationship to the legacy styles: the baseline block/arrow/group styles in
% style.tikzstyles is loaded by shared/preamble.tex / the standalone wrapper;
% this file ADDS the refined house language (flat pastel phase fills, thin-ink
% rounded boxes, dashed loop enclosures, elbow handoffs) and never redefines a
% baseline style name. New-style figures should use ONLY the ari* styles.
%
% inline-width: vendor-coloured override — gate-T2 marker: this is a shared
% STYLE file. minimum width / minimum height / text width are centralised
% here (exactly as in style.tikzstyles) and must NOT be re-specified inline
% in figure sources.

% --------------------------------------------------------------- palette
% Flat pastels on white. FILL ENCODES PHASE (which lifecycle phase owns the
% box); outlines, arrows, headings and body text are a single soft near-black
% ink; gray is reserved for loop enclosures, notes and secondary flows;
% `accent` is the only saturated color and is rationed to at most ONE
% semantic role per figure (e.g. the blocking verdict edge).
\definecolor{phaseidea}{HTML}{DBE5F6}%   lavender-blue — ideation / planning
\definecolor{phaseexp}{HTML}{FBE3C9}%    peach         — experimentation / execution
\definecolor{phasepaper}{HTML}{E9E9E9}%  light gray    — writing / publication
\definecolor{phasegate}{HTML}{D9EDDC}%   pale green    — gates / evaluation / verdicts
\definecolor{accent}{HTML}{D55E00}%      vermilion     — the ONE emphasis color (Okabe-Ito)
\definecolor{ariink}{HTML}{3C3C3C}%      soft near-black — outlines, arrows, headings
\definecolor{arimuted}{HTML}{8A8A8A}%    mid gray      — loop enclosures, notes, dashes

% --------------------------------------------------------------- styles
\tikzset{
  % Canvas rhythm — put on the picture itself: \begin{tikzpicture}[aricanvas].
  % Tight gaps INSIDE a phase column (0.5cm vertical), generous whitespace
  % BETWEEN phase columns (1.6cm horizontal): the phase-grouped-column layout.
  aricanvas/.style={node distance=0.5cm and 1.6cm,
                    line cap=round, line join=round},
  % ----- boxes: rounded rectangle, thin ink outline, flat pastel fill ------
  % Default fill is white; phase membership is declared with the fill, either
  % via the shorthands below or explicitly: \node[aribox, fill=phaseexp] ...
  aribox/.style={rectangle, rounded corners=3pt, draw=ariink,
                 line width=0.6pt, fill=white,
                 minimum width=2.6cm, minimum height=0.8cm, text width=2.4cm,
                 align=center, inner sep=3pt, font=\sffamily\footnotesize,
                 text=ariink},
  ariboxwide/.style={aribox, minimum width=3.6cm, text width=3.4cm},
  aribox idea/.style={aribox, fill=phaseidea},
  aribox exp/.style={aribox, fill=phaseexp},
  aribox paper/.style={aribox, fill=phasepaper},
  aribox gate/.style={aribox, fill=phasegate},
  % ----- phase heading: bold sans ink text above each column ---------------
  ariphase/.style={font=\bfseries\sffamily\small, text=ariink,
                   align=center, inner sep=2pt},
  % ----- loop enclosure: dashed gray rounded rectangle, fit-positioned -----
  %   \node[ariloop, fit=(a)(b)] (loop) {};
  ariloop/.style={rectangle, draw=arimuted, dashed, line width=0.7pt,
                  rounded corners=7pt, inner sep=8pt},
  % ----- arrows -------------------------------------------------------------
  % ariarrow        primary flow (short, solid, ink)
  % ariarrow muted  secondary / data / observation flow (dashed gray)
  % ariarrow accent the rationed emphasis edge (solid vermilion)
  % arifeedback     long outer feedback path (dotted gray, rounded routing)
  ariarrow/.style={-{Stealth[length=5pt,width=4pt]}, draw=ariink,
                   line width=0.7pt, shorten <=2.5pt, shorten >=2.5pt},
  ariarrow muted/.style={ariarrow, draw=arimuted, dashed},
  ariarrow accent/.style={ariarrow, draw=accent},
  arifeedback/.style={ariarrow, draw=arimuted, densely dotted,
                      rounded corners=7pt},
  % ----- elbow handoff between phase columns --------------------------------
  % Used with `to`: \draw[arielbow] (lastA) to (firstB);
  % arielbow    goes horizontal first, then vertical  (-|)
  % arielbow vh goes vertical first, then horizontal  (|-)
  arielbow/.style={ariarrow, rounded corners=5pt,
                   to path={-| (\tikztotarget) \tikztonodes}},
  arielbow vh/.style={ariarrow, rounded corners=5pt,
                      to path={|- (\tikztotarget) \tikztonodes}},
  % ----- small annotation / sparse edge label --------------------------------
  arinote/.style={font=\sffamily\scriptsize, text=arimuted, align=center,
                  fill=white, inner sep=1.5pt, rounded corners=1pt},
  % ----- numbered step badge (replaces edge labels for flow ordering) -------
  aribadge/.style={circle, fill=ariink, text=white, inner sep=1pt,
                   minimum size=11pt, font=\bfseries\sffamily\scriptsize},
  % ----- stroke-drawn heading glyphs (no icon packages) ----------------------
  % ~3mm tall, designed to sit left of an ariphase heading:
  %   \node[ariphase] (h) {\figlabel{Fxx.head.idea}};
  %   \pic at ([xshift=-2.5mm]h.west) {ariglyph bulb};
  ariglyph bulb/.pic={% ideation (lightbulb)
    \draw[ariink, line width=0.6pt] (0,0.45mm) circle [radius=0.95mm];
    \draw[ariink, line width=0.6pt] (-0.45mm,-0.85mm) -- (0.45mm,-0.85mm);
    \draw[ariink, line width=0.6pt] (-0.32mm,-1.3mm) -- (0.32mm,-1.3mm);
  },
  ariglyph flask/.pic={% experimentation (Erlenmeyer flask)
    \draw[ariink, line width=0.6pt]
      (-0.35mm,1.55mm) -- (-0.35mm,0.25mm) -- (-1.15mm,-1.45mm)
        -- (1.15mm,-1.45mm) -- (0.35mm,0.25mm) -- (0.35mm,1.55mm);
    \draw[ariink, line width=0.6pt] (-0.65mm,1.55mm) -- (0.65mm,1.55mm);
  },
  ariglyph doc/.pic={% writing / publication (page with dog-ear)
    \draw[ariink, line width=0.6pt]
      (-1.05mm,1.5mm) -- (-1.05mm,-1.5mm) -- (1.05mm,-1.5mm)
        -- (1.05mm,0.65mm) -- (0.2mm,1.5mm) -- cycle;
    \draw[ariink, line width=0.6pt] (-0.55mm,0.1mm) -- (0.55mm,0.1mm);
    \draw[ariink, line width=0.6pt] (-0.55mm,-0.6mm) -- (0.55mm,-0.6mm);
  },
  ariglyph gate/.pic={% gate / evaluation (check in circle)
    \draw[ariink, line width=0.6pt] (0,0) circle [radius=1.4mm];
    \draw[ariink, line width=0.6pt]
      (-0.65mm,0.05mm) -- (-0.15mm,-0.55mm) -- (0.75mm,0.6mm);
  },
  ariglyph loop/.pic={% iteration marker, sits inside an ariloop enclosure
    \draw[ariink, line width=0.6pt, -{Stealth[length=1.5pt,width=1.4pt]}]
      (115:1.2mm) arc[start angle=115, end angle=-50, radius=1.2mm];
    \draw[ariink, line width=0.6pt, -{Stealth[length=1.5pt,width=1.4pt]}]
      (-65:1.2mm) arc[start angle=-65, end angle=-230, radius=1.2mm];
  },
}%
}%
  {\IfFileExists{shared/figures/tikz/_style.tex}%
    {% shared/figures/tikz/_style.tex — ARI house design language for system figures.
% ---------------------------------------------------------------------------
% \input this file at the TOP of a figure source, BEFORE \begin{tikzpicture}.
% It only DEFINES colors / \tikzset styles / pics — it typesets nothing — so
% it is safe to \input from several figures in the same document and safe to
% load in any engine (pdflatex / lualatex / xelatex): xcolor + TikZ core only,
% no fontspec, no external icon packages (glyphs below are stroke-drawn pics).
%
% Relationship to the legacy styles: the baseline block/arrow/group styles in
% style.tikzstyles is loaded by shared/preamble.tex / the standalone wrapper;
% this file ADDS the refined house language (flat pastel phase fills, thin-ink
% rounded boxes, dashed loop enclosures, elbow handoffs) and never redefines a
% baseline style name. New-style figures should use ONLY the ari* styles.
%
% inline-width: vendor-coloured override — gate-T2 marker: this is a shared
% STYLE file. minimum width / minimum height / text width are centralised
% here (exactly as in style.tikzstyles) and must NOT be re-specified inline
% in figure sources.

% --------------------------------------------------------------- palette
% Flat pastels on white. FILL ENCODES PHASE (which lifecycle phase owns the
% box); outlines, arrows, headings and body text are a single soft near-black
% ink; gray is reserved for loop enclosures, notes and secondary flows;
% `accent` is the only saturated color and is rationed to at most ONE
% semantic role per figure (e.g. the blocking verdict edge).
\definecolor{phaseidea}{HTML}{DBE5F6}%   lavender-blue — ideation / planning
\definecolor{phaseexp}{HTML}{FBE3C9}%    peach         — experimentation / execution
\definecolor{phasepaper}{HTML}{E9E9E9}%  light gray    — writing / publication
\definecolor{phasegate}{HTML}{D9EDDC}%   pale green    — gates / evaluation / verdicts
\definecolor{accent}{HTML}{D55E00}%      vermilion     — the ONE emphasis color (Okabe-Ito)
\definecolor{ariink}{HTML}{3C3C3C}%      soft near-black — outlines, arrows, headings
\definecolor{arimuted}{HTML}{8A8A8A}%    mid gray      — loop enclosures, notes, dashes

% --------------------------------------------------------------- styles
\tikzset{
  % Canvas rhythm — put on the picture itself: \begin{tikzpicture}[aricanvas].
  % Tight gaps INSIDE a phase column (0.5cm vertical), generous whitespace
  % BETWEEN phase columns (1.6cm horizontal): the phase-grouped-column layout.
  aricanvas/.style={node distance=0.5cm and 1.6cm,
                    line cap=round, line join=round},
  % ----- boxes: rounded rectangle, thin ink outline, flat pastel fill ------
  % Default fill is white; phase membership is declared with the fill, either
  % via the shorthands below or explicitly: \node[aribox, fill=phaseexp] ...
  aribox/.style={rectangle, rounded corners=3pt, draw=ariink,
                 line width=0.6pt, fill=white,
                 minimum width=2.6cm, minimum height=0.8cm, text width=2.4cm,
                 align=center, inner sep=3pt, font=\sffamily\footnotesize,
                 text=ariink},
  ariboxwide/.style={aribox, minimum width=3.6cm, text width=3.4cm},
  aribox idea/.style={aribox, fill=phaseidea},
  aribox exp/.style={aribox, fill=phaseexp},
  aribox paper/.style={aribox, fill=phasepaper},
  aribox gate/.style={aribox, fill=phasegate},
  % ----- phase heading: bold sans ink text above each column ---------------
  ariphase/.style={font=\bfseries\sffamily\small, text=ariink,
                   align=center, inner sep=2pt},
  % ----- loop enclosure: dashed gray rounded rectangle, fit-positioned -----
  %   \node[ariloop, fit=(a)(b)] (loop) {};
  ariloop/.style={rectangle, draw=arimuted, dashed, line width=0.7pt,
                  rounded corners=7pt, inner sep=8pt},
  % ----- arrows -------------------------------------------------------------
  % ariarrow        primary flow (short, solid, ink)
  % ariarrow muted  secondary / data / observation flow (dashed gray)
  % ariarrow accent the rationed emphasis edge (solid vermilion)
  % arifeedback     long outer feedback path (dotted gray, rounded routing)
  ariarrow/.style={-{Stealth[length=5pt,width=4pt]}, draw=ariink,
                   line width=0.7pt, shorten <=2.5pt, shorten >=2.5pt},
  ariarrow muted/.style={ariarrow, draw=arimuted, dashed},
  ariarrow accent/.style={ariarrow, draw=accent},
  arifeedback/.style={ariarrow, draw=arimuted, densely dotted,
                      rounded corners=7pt},
  % ----- elbow handoff between phase columns --------------------------------
  % Used with `to`: \draw[arielbow] (lastA) to (firstB);
  % arielbow    goes horizontal first, then vertical  (-|)
  % arielbow vh goes vertical first, then horizontal  (|-)
  arielbow/.style={ariarrow, rounded corners=5pt,
                   to path={-| (\tikztotarget) \tikztonodes}},
  arielbow vh/.style={ariarrow, rounded corners=5pt,
                      to path={|- (\tikztotarget) \tikztonodes}},
  % ----- small annotation / sparse edge label --------------------------------
  arinote/.style={font=\sffamily\scriptsize, text=arimuted, align=center,
                  fill=white, inner sep=1.5pt, rounded corners=1pt},
  % ----- numbered step badge (replaces edge labels for flow ordering) -------
  aribadge/.style={circle, fill=ariink, text=white, inner sep=1pt,
                   minimum size=11pt, font=\bfseries\sffamily\scriptsize},
  % ----- stroke-drawn heading glyphs (no icon packages) ----------------------
  % ~3mm tall, designed to sit left of an ariphase heading:
  %   \node[ariphase] (h) {\figlabel{Fxx.head.idea}};
  %   \pic at ([xshift=-2.5mm]h.west) {ariglyph bulb};
  ariglyph bulb/.pic={% ideation (lightbulb)
    \draw[ariink, line width=0.6pt] (0,0.45mm) circle [radius=0.95mm];
    \draw[ariink, line width=0.6pt] (-0.45mm,-0.85mm) -- (0.45mm,-0.85mm);
    \draw[ariink, line width=0.6pt] (-0.32mm,-1.3mm) -- (0.32mm,-1.3mm);
  },
  ariglyph flask/.pic={% experimentation (Erlenmeyer flask)
    \draw[ariink, line width=0.6pt]
      (-0.35mm,1.55mm) -- (-0.35mm,0.25mm) -- (-1.15mm,-1.45mm)
        -- (1.15mm,-1.45mm) -- (0.35mm,0.25mm) -- (0.35mm,1.55mm);
    \draw[ariink, line width=0.6pt] (-0.65mm,1.55mm) -- (0.65mm,1.55mm);
  },
  ariglyph doc/.pic={% writing / publication (page with dog-ear)
    \draw[ariink, line width=0.6pt]
      (-1.05mm,1.5mm) -- (-1.05mm,-1.5mm) -- (1.05mm,-1.5mm)
        -- (1.05mm,0.65mm) -- (0.2mm,1.5mm) -- cycle;
    \draw[ariink, line width=0.6pt] (-0.55mm,0.1mm) -- (0.55mm,0.1mm);
    \draw[ariink, line width=0.6pt] (-0.55mm,-0.6mm) -- (0.55mm,-0.6mm);
  },
  ariglyph gate/.pic={% gate / evaluation (check in circle)
    \draw[ariink, line width=0.6pt] (0,0) circle [radius=1.4mm];
    \draw[ariink, line width=0.6pt]
      (-0.65mm,0.05mm) -- (-0.15mm,-0.55mm) -- (0.75mm,0.6mm);
  },
  ariglyph loop/.pic={% iteration marker, sits inside an ariloop enclosure
    \draw[ariink, line width=0.6pt, -{Stealth[length=1.5pt,width=1.4pt]}]
      (115:1.2mm) arc[start angle=115, end angle=-50, radius=1.2mm];
    \draw[ariink, line width=0.6pt, -{Stealth[length=1.5pt,width=1.4pt]}]
      (-65:1.2mm) arc[start angle=-65, end angle=-230, radius=1.2mm];
  },
}%
}%
    {% shared/figures/tikz/_style.tex — ARI house design language for system figures.
% ---------------------------------------------------------------------------
% \input this file at the TOP of a figure source, BEFORE \begin{tikzpicture}.
% It only DEFINES colors / \tikzset styles / pics — it typesets nothing — so
% it is safe to \input from several figures in the same document and safe to
% load in any engine (pdflatex / lualatex / xelatex): xcolor + TikZ core only,
% no fontspec, no external icon packages (glyphs below are stroke-drawn pics).
%
% Relationship to the legacy styles: the baseline block/arrow/group styles in
% style.tikzstyles is loaded by shared/preamble.tex / the standalone wrapper;
% this file ADDS the refined house language (flat pastel phase fills, thin-ink
% rounded boxes, dashed loop enclosures, elbow handoffs) and never redefines a
% baseline style name. New-style figures should use ONLY the ari* styles.
%
% inline-width: vendor-coloured override — gate-T2 marker: this is a shared
% STYLE file. minimum width / minimum height / text width are centralised
% here (exactly as in style.tikzstyles) and must NOT be re-specified inline
% in figure sources.

% --------------------------------------------------------------- palette
% Flat pastels on white. FILL ENCODES PHASE (which lifecycle phase owns the
% box); outlines, arrows, headings and body text are a single soft near-black
% ink; gray is reserved for loop enclosures, notes and secondary flows;
% `accent` is the only saturated color and is rationed to at most ONE
% semantic role per figure (e.g. the blocking verdict edge).
\definecolor{phaseidea}{HTML}{DBE5F6}%   lavender-blue — ideation / planning
\definecolor{phaseexp}{HTML}{FBE3C9}%    peach         — experimentation / execution
\definecolor{phasepaper}{HTML}{E9E9E9}%  light gray    — writing / publication
\definecolor{phasegate}{HTML}{D9EDDC}%   pale green    — gates / evaluation / verdicts
\definecolor{accent}{HTML}{D55E00}%      vermilion     — the ONE emphasis color (Okabe-Ito)
\definecolor{ariink}{HTML}{3C3C3C}%      soft near-black — outlines, arrows, headings
\definecolor{arimuted}{HTML}{8A8A8A}%    mid gray      — loop enclosures, notes, dashes

% --------------------------------------------------------------- styles
\tikzset{
  % Canvas rhythm — put on the picture itself: \begin{tikzpicture}[aricanvas].
  % Tight gaps INSIDE a phase column (0.5cm vertical), generous whitespace
  % BETWEEN phase columns (1.6cm horizontal): the phase-grouped-column layout.
  aricanvas/.style={node distance=0.5cm and 1.6cm,
                    line cap=round, line join=round},
  % ----- boxes: rounded rectangle, thin ink outline, flat pastel fill ------
  % Default fill is white; phase membership is declared with the fill, either
  % via the shorthands below or explicitly: \node[aribox, fill=phaseexp] ...
  aribox/.style={rectangle, rounded corners=3pt, draw=ariink,
                 line width=0.6pt, fill=white,
                 minimum width=2.6cm, minimum height=0.8cm, text width=2.4cm,
                 align=center, inner sep=3pt, font=\sffamily\footnotesize,
                 text=ariink},
  ariboxwide/.style={aribox, minimum width=3.6cm, text width=3.4cm},
  aribox idea/.style={aribox, fill=phaseidea},
  aribox exp/.style={aribox, fill=phaseexp},
  aribox paper/.style={aribox, fill=phasepaper},
  aribox gate/.style={aribox, fill=phasegate},
  % ----- phase heading: bold sans ink text above each column ---------------
  ariphase/.style={font=\bfseries\sffamily\small, text=ariink,
                   align=center, inner sep=2pt},
  % ----- loop enclosure: dashed gray rounded rectangle, fit-positioned -----
  %   \node[ariloop, fit=(a)(b)] (loop) {};
  ariloop/.style={rectangle, draw=arimuted, dashed, line width=0.7pt,
                  rounded corners=7pt, inner sep=8pt},
  % ----- arrows -------------------------------------------------------------
  % ariarrow        primary flow (short, solid, ink)
  % ariarrow muted  secondary / data / observation flow (dashed gray)
  % ariarrow accent the rationed emphasis edge (solid vermilion)
  % arifeedback     long outer feedback path (dotted gray, rounded routing)
  ariarrow/.style={-{Stealth[length=5pt,width=4pt]}, draw=ariink,
                   line width=0.7pt, shorten <=2.5pt, shorten >=2.5pt},
  ariarrow muted/.style={ariarrow, draw=arimuted, dashed},
  ariarrow accent/.style={ariarrow, draw=accent},
  arifeedback/.style={ariarrow, draw=arimuted, densely dotted,
                      rounded corners=7pt},
  % ----- elbow handoff between phase columns --------------------------------
  % Used with `to`: \draw[arielbow] (lastA) to (firstB);
  % arielbow    goes horizontal first, then vertical  (-|)
  % arielbow vh goes vertical first, then horizontal  (|-)
  arielbow/.style={ariarrow, rounded corners=5pt,
                   to path={-| (\tikztotarget) \tikztonodes}},
  arielbow vh/.style={ariarrow, rounded corners=5pt,
                      to path={|- (\tikztotarget) \tikztonodes}},
  % ----- small annotation / sparse edge label --------------------------------
  arinote/.style={font=\sffamily\scriptsize, text=arimuted, align=center,
                  fill=white, inner sep=1.5pt, rounded corners=1pt},
  % ----- numbered step badge (replaces edge labels for flow ordering) -------
  aribadge/.style={circle, fill=ariink, text=white, inner sep=1pt,
                   minimum size=11pt, font=\bfseries\sffamily\scriptsize},
  % ----- stroke-drawn heading glyphs (no icon packages) ----------------------
  % ~3mm tall, designed to sit left of an ariphase heading:
  %   \node[ariphase] (h) {\figlabel{Fxx.head.idea}};
  %   \pic at ([xshift=-2.5mm]h.west) {ariglyph bulb};
  ariglyph bulb/.pic={% ideation (lightbulb)
    \draw[ariink, line width=0.6pt] (0,0.45mm) circle [radius=0.95mm];
    \draw[ariink, line width=0.6pt] (-0.45mm,-0.85mm) -- (0.45mm,-0.85mm);
    \draw[ariink, line width=0.6pt] (-0.32mm,-1.3mm) -- (0.32mm,-1.3mm);
  },
  ariglyph flask/.pic={% experimentation (Erlenmeyer flask)
    \draw[ariink, line width=0.6pt]
      (-0.35mm,1.55mm) -- (-0.35mm,0.25mm) -- (-1.15mm,-1.45mm)
        -- (1.15mm,-1.45mm) -- (0.35mm,0.25mm) -- (0.35mm,1.55mm);
    \draw[ariink, line width=0.6pt] (-0.65mm,1.55mm) -- (0.65mm,1.55mm);
  },
  ariglyph doc/.pic={% writing / publication (page with dog-ear)
    \draw[ariink, line width=0.6pt]
      (-1.05mm,1.5mm) -- (-1.05mm,-1.5mm) -- (1.05mm,-1.5mm)
        -- (1.05mm,0.65mm) -- (0.2mm,1.5mm) -- cycle;
    \draw[ariink, line width=0.6pt] (-0.55mm,0.1mm) -- (0.55mm,0.1mm);
    \draw[ariink, line width=0.6pt] (-0.55mm,-0.6mm) -- (0.55mm,-0.6mm);
  },
  ariglyph gate/.pic={% gate / evaluation (check in circle)
    \draw[ariink, line width=0.6pt] (0,0) circle [radius=1.4mm];
    \draw[ariink, line width=0.6pt]
      (-0.65mm,0.05mm) -- (-0.15mm,-0.55mm) -- (0.75mm,0.6mm);
  },
  ariglyph loop/.pic={% iteration marker, sits inside an ariloop enclosure
    \draw[ariink, line width=0.6pt, -{Stealth[length=1.5pt,width=1.4pt]}]
      (115:1.2mm) arc[start angle=115, end angle=-50, radius=1.2mm];
    \draw[ariink, line width=0.6pt, -{Stealth[length=1.5pt,width=1.4pt]}]
      (-65:1.2mm) arc[start angle=-65, end angle=-230, radius=1.2mm];
  },
}%
}}%
\begin{tikzpicture}[aricanvas]

  % ---------- stage 1: rollout — paper + rubric in, submission out ------
  \node[ariphase]                              (h1)    {\figlabel{F16.rollout}};
  \node[aribadge] at ([xshift=-3.5mm]h1.west)          {1};
  \node[aribox, below=0.35cm of h1]            (paper) {\figlabel{F16.paper}};
  \node[ariboxwide, fill=phaseexp, below=of paper] (subm) {\figlabel{F16.submission}};
  \draw[ariarrow] (paper) -- (subm);

  % ---------- stage 2: reproduce — sandboxed execution of the submission
  \node[aribox exp, right=of subm]             (sandbox) {\figlabel{F16.sandbox}};
  \node[ariphase, above=0.35cm of sandbox]     (h2)      {\figlabel{F16.reproduce}};
  \node[aribadge] at ([xshift=-3.5mm]h2.west)            {2};
  \node[aribox, below=of sandbox]              (log)     {\figlabel{F16.log}};
  \draw[ariarrow] (subm) -- (sandbox);
  \draw[ariarrow] (sandbox) -- (log);

  % ---------- stage 3: grade — judge grades leaves, weights roll up -----
  \node[aribox, right=of log]                  (judge) {\figlabel{F16.judge}};
  \node[ariphase, above=0.35cm of judge]       (h3)    {\figlabel{F16.grade}};
  \node[aribadge] at ([xshift=-3.5mm]h3.west)          {3};
  \node[aribox gate, below=0.9cm of judge]     (l2)    {\figlabel{F16.pass}};
  \node[aribox gate, left=0.3cm of l2]         (l1)    {\figlabel{F16.pass}};
  \node[aribox, right=0.3cm of l2]             (l3)    {\figlabel{F16.fail}};
  \node[ariboxwide, fill=phasegate, below=0.9cm of l2] (score) {\figlabel{F16.score}};
  \draw[ariarrow] (log) -- (judge);
  \draw[ariarrow muted] (judge) -- (l1);
  \draw[ariarrow muted] (judge) -- node[arinote] {\figlabel{F16.grades}} (l2);
  \draw[ariarrow muted] (judge) -- (l3);
  \draw[ariarrow] (l1) -- (score);
  \draw[ariarrow] (l2) -- (score);
  \draw[ariarrow] (l3) -- node[arinote] {\figlabel{F16.weight}} (score);

\end{tikzpicture}
}
  \caption{The PaperBench replication path as three composable stages.
           The agent \emph{rollout} turns the paper and its rubric into a
           submission; \emph{reproduce} executes that submission in an
           isolated sandbox, capturing \texttt{reproduce.log} and its
           outputs; \emph{grade} scores the submission against the rubric
           tree. ARI's three-stage bridge orchestrates the stages and
           aggregates repeated runs into a single score.}
  \label{fig:pb-stages}
\end{figure*}

Two design choices keep this faithful to the benchmark. First, grading
is left to the vendored judge; ARI's adapters only reshape its output
and aggregate repeated runs into a single score, so upgrading
PaperBench is a vendor swap with no change to how leaf grades are
combined --- the same envelope feeds \cref{eq:paperscore}, with
$\judgefn{i}\in\{0,1\}$ for binary leaves. Second, the wrapper
\emph{fails loud}: when a benchmark-fair run cannot be served (no
container runtime, no scheduler, no GPU resource registration) it
raises rather than silently degrading to a weaker host, because a
silent downgrade would make the resulting score incomparable to a
leaderboard entry. Legacy silent-fallback behaviour is available only
behind an explicit opt-in.

A second concern is keeping the replication agent honest about what
the host actually provides. Left to the upstream prompt, an agent
tends to scaffold a reproduction script that calls binaries the host
lacks, and to default to Python regardless of the paper's real
language. ARI counters both by priming the agent to probe the
environment before it scaffolds, and by injecting a per-host notes
block --- generated from a live probe of the available toolchain, GPUs,
and network reachability --- so that what the agent reads matches what
the host can do.

\subsection{Domain breadth: the execution-profile contract}\label{sec:paper:paperbench:v072}

\begin{figure}[t]
  \centering
  \resizebox{\columnwidth}{!}{% F14: execution_profile contract (v0.7.2) — how the optional
% reproduce_contract.execution_profile field threads from a single
% rubric artefact into the BasicAgent prompt and the Phase 2 sbatch
% dispatcher, gated by three runtime probes that adapt the dispatched
% flag set to the local SLURM deployment.
% style.tikzstyles is loaded by shared/preamble.tex / standalone wrapper.
\begin{tikzpicture}[node distance=1.2cm and 0.45cm]

  % --- middle row: the two consumer drivers --------------------------
  \node[block native big]                                       (prompt)    {\figlabel{F14.prompt}};
  \node[block native big, right=4.2cm of prompt]                (sbatchn)   {\figlabel{F14.sbatch}};

  % --- top: the rubric contract sits above-centre --------------------
  \node[artifact wide, above=2.0cm of $(prompt)!0.5!(sbatchn)$] (rubric)    {\figlabel{F14.rubric}};

  % --- middle-left children (prompt blocks) --------------------------
  \node[block lane native, below=of prompt]                     (bshape)    {\figlabel{F14.b.shape}};
  \node[block lane native, left=of bshape]                      (bprof)     {\figlabel{F14.b.profile}};
  \node[block lane native, right=of bshape]                     (bconv)     {\figlabel{F14.b.conventions}};

  % --- middle-right children (sbatch probes) -------------------------
  \node[block lane native, below=of sbatchn]                    (probegres) {\figlabel{F14.probe.gres}};
  \node[block lane native, left=of probegres]                   (probefs)   {\figlabel{F14.probe.fs}};
  \node[block lane native, right=of probegres]                  (probehelp) {\figlabel{F14.probe.help}};

  % --- bottom: outcome -----------------------------------------------
  \node[artifact wide, below=1.4cm of $(bshape)!0.5!(probegres)$] (outcome) {\figlabel{F14.outcome}};

  % --- group bands ---------------------------------------------------
  \begin{scope}[on background layer]
    \node[group strong, fit=(prompt)(bprof)(bshape)(bconv)]              (bandL) {};
    \node[group strong, fit=(sbatchn)(probefs)(probegres)(probehelp)]    (bandR) {};
  \end{scope}
  \node[band label, above=4pt of bandL.north west, anchor=south west] {\figlabel{F14.band.prompt}};
  \node[band label, above=4pt of bandR.north west, anchor=south west] {\figlabel{F14.band.sbatch}};

  % --- edges ---------------------------------------------------------
  \draw[arrow] (rubric.south) to[bend right=15] (prompt.north);
  \draw[arrow] (rubric.south) to[bend left=15]  (sbatchn.north);

  \draw[arrow dashed] (prompt)  -- (bprof);
  \draw[arrow dashed] (prompt)  -- (bshape);
  \draw[arrow dashed] (prompt)  -- (bconv);

  \draw[arrow dashed] (sbatchn) -- (probefs);
  \draw[arrow dashed] (sbatchn) -- (probegres);
  \draw[arrow dashed] (sbatchn) -- (probehelp);

  \draw[arrow] (bshape)    -- (outcome);
  \draw[arrow] (probegres) -- (outcome);

\end{tikzpicture}
}
  \caption{The execution-profile contract threads from a single rubric
           artefact into two consumers: the replication agent's prompt
           (left) and the cluster dispatcher (right). Runtime probes
           adapt the dispatched resource request to the local scheduler,
           so one rubric runs faithfully on heterogeneous clusters
           without being modified.}
  \label{fig:execution-profile}
\end{figure}

PaperBench upstream assumes a single container with one GPU. To cover
papers whose methodology is fundamentally parallel --- multi-rank MPI
kernels, multi-node training, cluster-specific module environments ---
ARI extends the rubric with an optional \term{execution profile}
(\cref{fig:execution-profile}). The profile is domain-agnostic: it
names a \emph{parallel-execution shape} (single-CPU, single- or
multi-GPU, MPI, MPI+GPU) together with the paper's scale envelope and
any scheduler hints, rather than anything about the paper's scientific
field. When a paper is single-machine the profile is omitted and the
pipeline degenerates to the original single-node invocation: the
contract is additive, not breaking.

The profile feeds two consumers (\cref{fig:execution-profile}). It is
appended to the agent's prompt as a compact ``compute-node execution
conventions'' block --- shared-filesystem discipline, a preference for
the scheduler's launcher over a bare MPI launcher, environment
activation, and wall-clock wrapping --- and it parameterises the
cluster dispatch (node, GPU, memory, and placement requests). Because
clusters disagree on which scheduler flags they accept, the dispatcher
probes the local scheduler at run time and drops flags it does not
support, so the same rubric dispatches faithfully on a multi-GPU
cluster, on a sandbox partition without GPU registration, or on a
single workstation. A regression guard enforces the inverse: ARI
running inside an unrelated cluster allocation must never leak
cluster- or MPI-specific prose into the prompt for a non-HPC paper.
Only this wrapper layer changes; the vendored benchmark and ARI core
are untouched across the extension.

\subsection{Failure modes and guardrails}\label{sec:paper:failures}

Three failure modes are common enough to be named, each with a
structural --- not merely advisory --- mitigation:

\begin{itemize}
  \item \textbf{LLM-only support}: the paper asserts a claim that no
        node in the lineage substantiates. The mitigation constrains
        the writer to a pre-selected set of lineage artefacts, so every
        numerical or causal claim is traceable to a node that produced
        it; the draft cannot cite evidence that was not generated.
  \item \textbf{Citation hallucination}: the paper cites a non-existent
        work. The mitigation is structural --- the writer may not invent
        citation keys, and every bibliography entry is verified against
        an external index before it is admitted (\cref{sec:related}).
  \item \textbf{Refine drift}: the aggregate score climbs while the
        paper loses coherence. The mitigation is a per-axis floor: if
        any axis falls below a threshold the refine is aborted, so a
        gain on one axis cannot mask a collapse on another.
\end{itemize}

The first guard also bounds context: the writer is given only the
bytes of the selected lineage artefacts, under a size budget, so the
draft stays within the model's context window even for long runs.
